% Imporditud: tools/import_eksperimendid.py
% Allikas: Eesti füüsikaolümpiaadi eksperimendi ülesannete
%   kogu (Rael Kalda, 2023), ülesanne Ü98, lahendus L98.
\setAuthor{EFO žürii}
\setRound{piirkonnavoor}
\setYear{2008}
\setNumber{G E1}
\setDifficulty{3}
\setTopic{vedelike mehaanika}
\setVahendid{joonlaud, anum nr. 1 tuntud vedelikuga (tihedus $\rho_1$ = 1000 kg/m$^3$), anum nr. 2 tuntud vedelikuga (tihedus $\rho_2$ = 1060 kg/m$^3$), anum tundmatu vedeli-}
\setPunktid{10}
\setAllikas{Eksperimendi kogu Ü98, lahendus lk 76}
\setLahendusseis{toores}

\prob{Vedeliku tihedus}
Määrata tundmatu vedeliku tihedus. Hinnata mõõteviga.

kuga, 2 joogikõrt, käärid, plastiliin.

\hint

\solu
Valmistame kõrrest ja plastiliinist tiheduse mõõtmiseks vajaliku seadme --- areomeetri (1 p). Olgu areomeetri kogumass m ning sukeldunud ruumalad vastavalt $V_1$ (tuntud vedelik nr. 1), $V_2$ (tuntud vedelik nr. 2) ja $V_3$ (tundmatu vedelik nr. 3). Areomeetri vedelikus ujumise korral on täidetud tingimus

$\rho$1V1 = $\rho$2V2 = $\rho$3V3 = m (1 p).

Võtame kõrrel (ristlõikepindalaga S ) null-nivooks sukeldumissügavuse vedeliku ``3'' puhul ning olgu sukeldumissügavused ülejäänud kahe vedeliku korral vastavalt $h_1$ ja $h_2$ (need on märgiga suurused). Siis

$V_1$ = $V_3$ + h1S, $V_2$ = $V_3$ + h2S (1 p).

Kirjutame esimesed võrrandid ümber kujul

$V_3$ = m m m , $V_3$ + Sh2 = , $V_3$ + Sh1 = $\rho_1$ $\Rightarrow \rho$3 $\rho_2$

1 $-$ 1 = Sh2 , 1 $-$ 1 = Sh1 (1 p). $\rho_2$ $\rho_3$ m $\rho_1$ $\rho_3$ m

Sellest võrrandisüsteemist saame avaldada otsitava tiheduse $\rho_3$ :

$\rho_3$ = $h_1$ $-$ $h_2$ (1 p). $h_1$ $-$ $h_2$ $\rho_2$ $\rho_1$

Mõõtmised: Mõõtmised on dokumenteeritud ja nende põhjal on leitud $\rho$ (1 p). tulemus on tõepärane (1 p). Mõõtmisi on korratud mitu korda (1 p). Hinnatud on mõõteviga (2 p).

Alternatiivne (ebatäpne) lahendus: Eeldame, et areomeetri sukeldumissügavus on võrdeline vedeliku tihedusega:

$h_1$ = k ($\rho_1$ $- \rho$3) , $h_2$ = k ($\rho_2$ $- \rho$3) (1 p),

kus nivood on jällegi mõõdetud tundmatu vedeliku sukeldumissügavuse suhtes.

Siit leiame $\rho$2h1 $-$ $\rho$1h2 $h_1$ $-$ $h_2$ $\rho_3$ = (1 p).

Niisiis saab sellise lahenduse puhul lõppvalemi eest 4 p asemel 2 p. Juhul, kui seesama ligikaudne valem on tuletatud põhialustest (Archimedese seadus) lähtudes, kuid on kasutatud lähendust |$\rho_1$ $- \rho$2| $\ll \rho$1 ja |$\rho_1$ $- \rho$3| $\ll \rho$1, siis lisandub üks punkt (st valemi eest kokku 3 p).
\probend
